\chapter{Referencial Teórico}
\label{cap-referencial-teorico}

% Instruções gerais para o Referencial Teórico:
% 1. Este capítulo deve apresentar a fundamentação teórica do seu trabalho
% 2. Organize os conceitos do mais amplo para o mais específico
% 3. Use citações para apoiar seus argumentos
% 4. Crie seções lógicas e conectadas
% 5. Mantenha o foco nos temas relevantes para seu trabalho

\section{Primeiro Conceito Principal}
\label{rt-conceito1}
% Desenvolva aqui o primeiro conceito importante para seu trabalho
% Use subseções se necessário para organizar melhor o conteúdo

% ------ DENGE | EPIDEMIOLOGIA | BRASIL -----------
% CICLO DE VIDA DO MOSQUITO E TRANSMISSÃO DO VÍRUS:
% Após se alimentar de um hospedeiro infectado, a fêmea do mosquito passa a ser vetora do vírus pelo restante de sua vida, uma vez que o agente patogênico replica-se em seu intestino e migra para as glândulas salivares. O intervalo entre a ingestão do vírus pelo mosquito e sua capacidade de transmiti-lo a um novo hospedeiro é denominado Período de Incubação Extrínseca (PIE). Esse período dura aproximadamente de 8 a 12 dias em temperaturas entre 25 °C e 28 °C, podendo variar conforme a amplitude das flutuações térmicas diárias, o genótipo viral e a concentração viral inicial \cite{who2024dengue}.

% CEÁRIO EPIDEMIOLÓGICO CRÍTICO:
% Usar o **Risk Factor Identification and Spatiotemporal Diffusion Path During the Dengue Outbreak** para falar como as mudanças estão afetando "novas áreas"

% CONTROLE VETORIAL vs. VACINAÇÃO:

%Em resposta a esse desafio, o Brasil tornou-se o primeiro país do mundo a disponibilizar vacinas contra a dengue no sistema público de saúde. A vacina Qdenga, indicada para indivíduos que ainda não contraíram a doença, demonstrou ser eficaz contra o DENV-1 em 69,8\% dos casos, contra o DENV-2 em 95,1\% e contra o DENV-3 em 48,9\%, com eficácia geral comprovada de 84,1\% na prevenção de hospitalizações \cite{agenciagov2024qdenga}. Embora a incorporação do imunizante ao Sistema Único de Saúde (SUS) em 2023 e ao Calendário Nacional de Vacinação em 2024 represente um avanço inquestionável, alcançando 521 dos 5.569 municípios brasileiros em sua primeira campanhas, o controle do vetor \textit{Aedes aegypti} permanece como principal estratégia de prevenção e controle da doença e de outras arboviroses urbanas como chikungunya e Zika \cite{ministerio2026dengue}.
% A complexidade do controle vetorial, contudo, é acentuada por um aspecto estrutural da dinâmica epidêmica: quando a transmissão já está estabelecida em uma região, com um ou mais sorotipos em circulação e elevada densidade de mosquitos, as ações reativas de controle mostram-se pouco efetivas. A velocidade de circulação viral supera amplamente o tempo necessário para a redução das populações do vetor. Nesses cenários, a queda no número de casos tende a ocorrer mais em função da imunidade de grupo adquirida pela população do que pela eficácia das intervenções aplicadas \cite{ministerio2026dengue}. Tal constatação reforça a importância crítica da atuação preventiva, tornando essencial o desenvolvimento de novas metodologias que apoiem a gestão pública antes da instauração de uma epidemia.

% Exemplo de como usar uma citação:
%\begin{citacao}
%A engenharia de software tem evoluído significativamente nos últimos anos \cite{inmetro2003}.
%\end{citacao}

%\begin{citacao}
%O desenvolvimento de modelos para turbulência tem sido um desafio constante na engenharia \cite{bordalo1989}.
%\end{citacao}

\section{Segundo Conceito Principal}
\label{rt-conceito2}

% ------ MULTIAGENTES -----------
% Desenvolva aqui o segundo conceito importante
% Estabeleça conexões com o conceito anterior quando pertinente

\section{Terceiro Conceito Principal}
\label{rt-conceito3}

% ------ DADOS -----------
% Continue desenvolvendo os conceitos relevantes
% Mantenha o foco na relevância para seu trabalho

% TODO - Mover para referencial:
% o Índice de Breteau e a densidade populacional são os fatores dominantes durante o surto ativo, e as variáveis meteorológicas apresentam correlação fraca ou não significativa em janelas curtas
% O Índice de Moran, também aplicado por Li e Chen (2016), pode ser usado para medir a autocorrelação espacial dos casos gerados pela simulação, desde que dados equivalentes estejam disponíveis para o cenário escolhido.


\section{Resumo do Capítulo}
\label{rt-resumo}
% Sintetize os principais conceitos discutidos
% Estabeleça as conexões entre os diferentes tópicos abordados